\newpage
\section{Benchmark Suite}\label{sec:benchmark-suite}

As described in the evaluation criteria~\ref{subsec:criteria} there will be at least 120 evaluation runs,
which will be a lot of data to analyze and assess.
For this reason, we decided to implement a benchmarking suite, consisting of a benchmark harness that executes benchmark runs and summarizes their results into a report.
The idea is to fully automate the aggregation of the \textbf{quantitative criteria} (duration, token usage) and
to prepare the reports needed for the \textbf{qualitative criteria} (correctness and consistency) as good as possible.
The key concepts for this implementation are described the table below~\ref{tab:benchmark-suite-concepts}:

\begin{table}[h]
    \centering
    \begin{tabular}{ll}
        \hline
        \textbf{Term} & \textbf{Concept} \\
        \hline
        \mintinline{python}{BenchmarkManager} & The automation harness \\
        \mintinline{python}{BenchmarkRun} & Represents a single execution \\
        \mintinline{python}{BenchmarkResult} & The result data type of a single execution \\
        \mintinline{python}{BenchmarkReport} & The collected output aggregating all results \\
        \hline
    \end{tabular}
    \caption{Benchmark Suite Concepts}
    \label{tab:benchmark-suite-concepts}
\end{table}
